% !TeX root = closed-systems-from-comparison-completeness.tex
% ============================================================
\chapter{Minimal Examples and Failure Modes}
\label{app:minimal-examples}
% ============================================================
\begin{chapterorientation}
\orientationitem{Purpose}{This appendix supplies small models that test the definitions used in the main text.}
\orientationitem{Status}{The examples are illustrative.  They do not replace the proofs in the chapters where the corresponding results are established.}
\orientationitem{Use}{Each example identifies the relevant theorem or definition and the failure mode it is meant to isolate.}
\end{chapterorientation}

The body of the monograph is theorem-first.  This appendix records
minimal examples and failure modes so that the main definitions can be
checked without carrying the full stack at once.  Each example is
deliberately small.  Its role is to isolate one structural distinction,
not to model a physical system in full.

\section{A missing rectangle}
\label{app:example-missing-rectangle}
Let
\[
	A=\{a_0,a_1\},\qquad B=\{b_0,b_1\},
\]
and consider the subset
\[
	U=\{(a_0,b_0),(a_1,b_0),(a_0,b_1)\}\subset A\times B.
\]
The left and right coordinates define two profile maps.  The realized
left profiles are \(a_0,a_1\), and the realized right profiles are
\(b_0,b_1\), but the pair \((a_1,b_1)\) is absent.

Thus \(U\) has the two marginal profile families but not their full
product.  In the language of \cref{def:rect}, it is not rectangularly
complete.  Adding the missing state \((a_1,b_1)\) gives the completion
\[
	U^{+}=A\times B.
\]
No old coordinate pair is changed by this addition.  The example is the
finite prototype for the distinction used in
\cref{lem:conservative-completion-dichotomy}: an omitted profile is
either internally excluded by the comparison predicates or it is a
conservative omission.

\section{A conservative omission is not closed}
\label{app:example-conservative-omission}
Continue with \(U\subset A\times B\) from
\cref{app:example-missing-rectangle}.  Suppose the comparison predicates
only read the left profile, the right profile, and equality of those
profiles where they are defined.  Then adding \((a_1,b_1)\) does not
alter the comparison values among the old states.  The omission is
therefore conservative relative to the old comparison algebra.

The point is not that every omission is conservative.  The point is that
a closed intrinsic comparison world cannot omit a profile whose
realization would preserve all old finite-coordinate comparison data.
This is the finite example behind
\cref{cor:closure-forbids-conservative-omission} and the
rectangular-completeness theorem
\cref{thm:rectangular-completeness-intrinsicality}.

\section{A quotient report and a non-descending report}
\label{app:example-nondescending-report}
Let \(X=\{x_0,x_1\}\), and let \(G=\mathbb Z/2\mathbb Z\) act by
interchanging \(x_0\) and \(x_1\).  The quotient
\[
	\pi:X\to X/G
\]
has a single orbit.  A report
\[
	R:X\to\{0,1\},\qquad R(x_0)=0,\quad R(x_1)=1,
\]
does not descend to the quotient, because it is not constant on the
orbit fibre.  The constant report \(R_0(x_0)=R_0(x_1)=0\) does descend.

This is the minimal form of the descent condition in
\cref{thm:descent}: coherent physical reports are precisely those
representative-level reports that are constant on the quotient fibres.
Representative labels may still be useful in a calculation, but they
are not physical data unless the report descends.

\section{Twin supported descriptions}
\label{app:example-twin-descriptions}
Let \(\Gamma(t)\) and \(\widetilde\Gamma(t)\) be two representative
paths in a \(G\)-space \(X\), with
\[
	\widetilde\Gamma(t)=g(t)^{-1}\cdot\Gamma(t).
\]
Then
\[
	\pi(\Gamma(t))=\pi(\widetilde\Gamma(t))
\]
for every \(t\).  A description in which the \(A\)-side is held fixed
and the \(B\)-side carries the motion, and a twin description in which
the \(B\)-side is held fixed and the \(A\)-side carries the motion, can
therefore determine the same quotient path in \(\Phys=X/G\).

No coherent report distinguishes the two descriptions unless it uses
representative data that fail to descend.  This is the small model
behind the orbit-equivalent supported representatives of
\cref{ch:rotation}.  It records the logical content of the twin-loop
figure: the equality is equality of physical report values, not equality
of representative states or comparison arrows.

\section{Endpoint data do not determine route data}
\label{app:example-route-fibre}
Let a protocol category contain two distinct paths
\[
	p,q:a\to b
\]
with the same source and target.  Let a route assignment \(H\) satisfy
\[
	H(p)\ne H(q).
\]
Then \(H\) cannot be a function of endpoint data alone.  It is not
constant on the endpoint fibre.  Conversely, any assignment that is
constant on every endpoint fibre factors through endpoint data.

This is the elementary form of the factorization criterion used
throughout the transport chapters.  Holonomy, loop defect, and
descent obstruction enter exactly when route data carry information not
recoverable from endpoints alone.

\section{A ghost carrier}
\label{app:example-ghost-carrier}
Let \(K\) be a two-dimensional real vector space with basis
\(\{e_1,e_2\}\), and let a realized second-jet comparison map be
\[
	J:K\to\mathbb R,
	\qquad
	J(e_1)=1,\quad J(e_2)=0.
\]
Then \(e_2\ne 0\) is an intrinsic carrier direction that is invisible to
the realized second-jet reading.  It lies in the ghost kernel of \(J\).

This is the finite linear prototype for the ghost-sector analysis in
\cref{sec:interface-ghosts}.  Condition \textup{(D)} rules out precisely
this kind of invisible intrinsic carrier in the realized branch, while
condition \textup{(T)} controls the torsion-side ambiguity.  Without the
ghost-free regime, the smooth realization may fail to read all
intrinsic degree--\(2\) obstruction data.

\section{Local flatness and global closure}
\label{app:example-local-flat-global-closed}
Let \(S^1_R\) be a circle of radius \(R\), parameterized by arc length.
For separations \(d\ll R\), an interval in \(S^1_R\) is locally
indistinguishable from an interval in \(\mathbb R\) up to curvature
corrections of higher order.  Globally, however, the circle closes:
moving far enough along the same geodesic returns to the starting
point.

The higher-dimensional spherical case has the same logical form.  Local
coordinates may read a small region as approximately flat, while the
global topology is closed.  This example is not a proof of the
manifold-stage \(S^3\) result.  It records the distinction needed in
\cref{ch:connection}: local smooth realization and global closure are
different questions.

\section{A normalization obstruction}
\label{app:example-normalization}
Suppose a positive scalar channel determines only ratios
\[
	\frac{s_i}{s_j}.
\]
Then multiplying every scalar by the same positive constant leaves all
ratios unchanged.  No internal comparison among ratios can select the
absolute unit.  An external calibration would fix the scale, but that
calibration is not produced by the ratio data themselves.

This is the minimal form of the normalization obstruction used in the
numerical closure chapters.  Internal structure can determine exact
dimensionless ratios while leaving an absolute normalization as boundary
data.

\begin{chapterclosing}
\closingitem{Established}{The principal definitions have been tested against finite examples and failure modes.}
\closingitem{Carried forward}{The examples provide a compact check on the abstract definitions in the main theorem chain.}
\end{chapterclosing}
